\documentclass{article}
\usepackage{amsmath, amssymb, amsthm}
\usepackage[margin=1in]{geometry}

\newcommand{\E}{\mathbb{E}}
\newcommand{\Prob}{\mathbb{P}}

\newtheorem{definition}{Definition}
\newtheorem{observation}{Observation}

\title{The Early-Arrival Strategy}
\author{}
\date{}

\begin{document}
\maketitle

\section{Motivation}

The original wait-time strategy (``late arrival'') was designed to position the particle at the end of the ON window, with the intuition that this would maximize future flexibility. However, simulation results revealed that a simpler ``early arrival'' strategy significantly outperforms it, particularly when the lap time $\tau$ is small relative to the ON period $T_1$.

This document describes the early-arrival strategy, explains why it works, and identifies the regime where it excels.

\section{Setup}

Recall the problem parameters:
\begin{itemize}
    \item Particle velocity $v > 0$, giving lap time $\tau = 2/v$
    \item Lifespan $l \sim \text{Exp}(\lambda)$, refreshed after each successful lap
    \item Gate ON during $(0, T_1]$, OFF during $(T_1, T_1 + T_2]$, with period $T = T_1 + T_2$
    \item Current offset $\Omega \in (0, T_1]$: position in the gate cycle when a lap begins
\end{itemize}

A particle at offset $\Omega$ with lifespan $l$ must choose a wait time $w \geq 0$. After waiting $w$ and traveling for $\tau$, it arrives at:
\[
\Omega' = (\Omega + w + \tau) \mod T
\]
The lap succeeds if $\Omega' \in (0, T_1]$ and $w + \tau \leq l$.

\section{The Early-Arrival Strategy}

\begin{definition}[Early-Arrival Strategy]
Given state $(\Omega, l)$, the early-arrival strategy chooses the minimum wait time $w^*$ such that:
\begin{enumerate}
    \item The arrival $\Omega' = (\Omega + w^* + \tau) \mod T$ lies in the ON window $(0, T_1]$
    \item The particle survives: $w^* + \tau \leq l$
\end{enumerate}
If no such $w^*$ exists, the particle cannot complete a lap and dies.
\end{definition}

Concretely:
\begin{itemize}
    \item If $(\Omega + \tau) \mod T \in (0, T_1]$, then $w^* = 0$ (immediate departure lands in ON window)
    \item Otherwise, compute the minimum wait to reach the next ON window
    \item If even this minimum wait exceeds available lifespan, fall back to the late-arrival strategy (which targets the latest reachable ON window)
\end{itemize}

\section{Why Early Arrival Works}

\subsection{The ``Spam Laps'' Phenomenon}

When $\tau < T_1$, a particle can complete multiple laps within a single ON period. Consider a particle at offset $\Omega$ with $\Omega + \tau \leq T_1$:
\begin{itemize}
    \item Immediate departure: arrive at $\Omega' = \Omega + \tau$, still in ON window
    \item Get fresh lifespan $l'$
    \item If $\Omega' + \tau \leq T_1$, can immediately do another lap
    \item Repeat until either $\Omega^{(k)} + \tau > T_1$ or unlucky lifespan
\end{itemize}

Each lap takes only $\tau$ time but grants a completely fresh lifespan draw. This creates a ``re-rolling'' effect: short lifespans are quickly replaced, while a lucky long lifespan provides a buffer for future laps.

\subsection{Phase Preservation}

The early strategy preserves the particle's phase relationship with the gate:
\[
\Omega' = \Omega + \tau \pmod{T}
\]
If $\Omega$ is in a favorable position (where immediate departure lands in ON), then $\Omega'$ is often also favorable. The particle stays synchronized with the gate.

In contrast, the late-arrival strategy actively repositions to $\Omega' = T_1$ every lap, which may not be optimal.

\subsection{Value Function Structure}

Simulation reveals that $U(\Omega)$---the expected laps starting from offset $\Omega$---is \emph{increasing} in $\Omega$ within certain regions. This is counterintuitive: being closer to the end of the ON window (larger $\Omega$) yields higher expected laps.

The explanation: larger $\Omega$ means the next OFF period is closer, but after crossing it, the particle emerges early in the next ON period with maximum runway. The value $U$ balances:
\begin{itemize}
    \item Immediate survival probability (favors small $\Omega$, more room in current window)
    \item Future positioning (favors large $\Omega$, better phase after crossing OFF)
\end{itemize}

For fast particles ($\tau \ll T_1$), immediate survival is nearly guaranteed for any $\Omega$, so future positioning dominates.

\section{Regime Analysis}

\begin{observation}[Fast Regime: $\tau < T_1$]
Early arrival significantly outperforms late arrival. Simulation shows improvements of 40--60\% in expected laps. The particle can ``spam laps'' to re-roll lifespans rapidly.
\end{observation}

\begin{observation}[Medium Regime: $T_1 < \tau < T$]
Early and late strategies perform similarly. Each lap necessarily crosses into the OFF period, so phase positioning matters less.
\end{observation}

\begin{observation}[Slow Regime: $\tau > T$]
Both strategies reduce to similar behavior---the particle wraps around multiple gate cycles per lap. Strategy differences are minimal.
\end{observation}

\section{Comparison with Upper and Lower Bounds}

For a particle with velocity $v$ and mean lifespan $1/\lambda$:

\textbf{Upper bound} (no gate): 
\[
\E[\text{laps}] = \frac{e^{-\lambda\tau}}{1 - e^{-\lambda\tau}} = \frac{1}{e^{\lambda\tau} - 1}
\]

\textbf{Lower bound} (gate, no waiting): Particles depart immediately regardless of phase. Many die by arriving during OFF.

\textbf{Early strategy}: Achieves values between these bounds, closer to upper bound when $\tau < T_1$.

\textbf{Late strategy}: Also between bounds, but typically 30--50\% lower than early in the fast regime.

\section{Open Questions}

\begin{enumerate}
    \item \textbf{Optimality}: Is early-arrival actually optimal, or is there a better strategy? The MDP framework suggests the optimal strategy maximizes $U(\Omega')$ over reachable $\Omega'$. If $U$ is monotonic, either pure-early or pure-late is optimal. But $U$ may have more complex structure.
    
    \item \textbf{Analytical form}: Can we derive $U(\Omega)$ analytically for the early strategy? The fixed-point equation is:
    \[
    U(\Omega) = \int_0^\infty \lambda e^{-\lambda l} \cdot V(\Omega, l) \, dl
    \]
    where $V(\Omega, l) = 1 + U(\Omega')$ if a lap is possible, else 0.
    
    \item \textbf{Threshold policies}: Perhaps the optimal strategy is neither pure-early nor pure-late, but a threshold policy: ``go early if $l < L^*(\Omega)$, else position carefully.'' This would balance re-rolling (for short $l$) against strategic positioning (for long $l$).
\end{enumerate}

\section{Conclusion}

The early-arrival strategy succeeds by maximizing the lap rate. Each lap is cheap (costs only $\tau$ time) but valuable (refreshes lifespan entirely). In the fast regime, this ``spam laps'' approach dramatically outperforms careful positioning.

The key insight: when survival per lap is nearly certain, \emph{more laps is better than better-positioned laps}.

\end{document}
